\section*{Acknowledgments}
The author used AI assistants in this work. OpenAI ChatGPT (\mbox{GPT-5.6 Luna}) and Google Gemini (\mbox{Gemini 3.8 Flash}) were each given the same draft for independent critical reading and cross-checking; this cross-check pointed out that the strips do not vanish outside the fundamental domain, which led to the moment conditions (H1) and (H2). Anthropic Claude (\mbox{claude-opus-5-5}) was used in two separate sessions. In the first session it took part in theoretical discussion, derived the gap formula (Proposition~\ref{prop:formula}) independently, wrote an independent evaluator of the strips, proposed the weaker target of Proposition~\ref{prop:target}, checked the proofs of the main results and reproduced key values (the strips and their valleys) by independent computations, and drafted the text. In the second session it designed and ran the numerical study, found the strips and the other families, derived Theorem~\ref{thm:direct}, Lemmas~\ref{lem:smoothing1d}, \ref{lem:heisenberg}, \ref{lem:split}, \ref{lem:fisher} and~\ref{lem:dip}, Theorems~\ref{thm:log}, \ref{thm:valley} and~\ref{thm:product} and Proposition~\ref{prop:novalley}, and wrote the evaluation and verification code. Results were passed between the sessions as written summaries and checked in the first session by re-running the verification suite and by independent methods. The author directed the research, made all decisions on its content, and takes full responsibility for this article.

\section*{Data and Code Availability}
Scripts, results and the verification suite that reproduce every number and figure of this article are openly available at \url{https://github.com/neginyan/resolution-gap-log-law}, together with independent checks and a notebook that reproduces the main numbers in a browser. The version used for this article is archived at Zenodo, \pending{code DOI}.
